近年、AIの発展は著しく、物理学や数学をはじめとする自然科学の分野で広く活用されつつあります。
一方で、人間社会の現象を分析する社会科学においては、AIの活用はいまだ発展途上にあります。

\textbf{経済学}は、社会科学の中でも大きな影響力を持つ学問です。
人類社会をより豊かにするため、経済現象を解明してきました。
そこでは主に、理論的なアプローチが用いられてきました。
人間の行動に対する数学的な仮定を置き、演繹的に結果を導く、というアプローチです。

経済現象をより深く理解するためには、理論だけでなく、実験による検証も重要です。
ところが、人間社会そのものを実験対象とすることには、倫理的制約やコストの問題があります。
そのため、経済学はしばしば「実験が難しい学問」とされてきました。

私たちは、この限界を乗り越えるために、\textbf{経済現象を計算機上でシミュレーション}することを目指しています。
大規模言語モデル(LLM)を用いて個々の人間や企業の行動をモデル化し、それらが相互作用することで生じる経済現象を分析します。
具体的には、100体のLLMによって消費者や企業をモデル化し、私たちが構築した経済理論に基づく動学環境の中で、相互作用させました。
その結果、「経済成長」や「国際貿易」といった既知の経済現象が再現されることを確認しました。

さらに私たちは、「文明崩壊級の隕石の接近」というシナリオも実験しました。
その結果、人々は愛する人と最期の時間を過ごすために労働を放棄し、企業は投資や研究開発といった長期的戦略を捨て去りました。
結果として、経済全体が崩壊していく様子が観察されました。

このように、現実社会では実験が困難な状況を検証できる点に、計算機シミュレーションの大きな価値があります。
今後は、より多様な経済現象を分析できるようシミュレーション環境を精緻化します。
「LLMが暮らす計算機上の経済社会」の構築を目指して研究を進めていきます。
